% TEMPLATE-MILC-v3.tex - plantilla maestra UNICA de las guias MILC v3.
%
% La usan las tres familias de guias, cada una con su driver:
%   · Grados 6-11  -> scripts/build-guias-g11.py
%   · Bebras       -> scripts/build-guias-bebras.py
%   · Semillero    -> scripts/build-guias-semillero.py
%
% Antes habia tres copias casi identicas (~400 lineas iguales, 6 distintas):
% cada mejora habia que aplicarla tres veces, y en la practica no se hacia
% (el motor de recursos vivia solo en dos de ellas). Lo que varia entre
% familias esta parametrizado en 6 placeholders que cada driver llena
% (se nombran aqui SIN su sintaxis real: el builder reemplaza por texto plano,
% tambien dentro de los comentarios, y un valor multilinea rompe el preambulo):
%
%   HEADER_IZQ        encabezado izquierdo de pagina
%   HEADER_DER        encabezado derecho de pagina
%   PORTADA_ETIQUETA  rotulo sobre el numero grande (GRADO/MOMENTO/MODULO)
%   PORTADA_SERIE     linea de serie de la portada
%   PORTADA_ID        identificador de la guia en la portada
%   PORTADA_CIRCULO   el circulito de grado (°), o vacio si no aplica
%
% El resto de claves las documenta content/guias/_SCHEMA.md. El script valida
% que no queden placeholders sin reemplazar antes de escribir el archivo final.

\documentclass[11pt,letterpaper]{article}
\usepackage[margin=1.75cm]{geometry}
\usepackage[table]{xcolor}
\usepackage{fontspec}
\usepackage[spanish]{babel}
\usepackage{tikz}
% Librerías TikZ para los diagramas de las guías (bloque `recursos:`):
%  · babel  → babel en español vuelve activos caracteres como «>», y TikZ los usa
%             en sus opciones (>=stealth, ->). Sin esta librería, cualquier
%             diagrama con flechas revienta con «\language@active@arg> has an
%             extra }». Debe ir ANTES de que se dibuje nada.
%  · positioning        → sintaxis «below left=2pt and 3pt».
%  · decorations.pathreplacing → llaves ({brace}) para señalar rangos.
\usetikzlibrary{babel,positioning,decorations.pathreplacing}
\usepackage{graphicx}
% Circuitos: el trabajo de electrónica se hace hoy con micro:bit y sus sensores
% integrados (MakeCode, sin cableado), así que no se necesitan esquemáticos.
% Si algún día entran montajes con protoboard, basta descomentar esta línea:
% los diagramas .tex/.tikz declarados en `recursos:` ya se incrustan solos.
% \usepackage{circuitikz}
\usepackage[most]{tcolorbox}
\usepackage{tabularx}
\usepackage{array}
\usepackage{fancyhdr}
\usepackage{titlesec}
\usepackage{ragged2e}
\usepackage{enumitem}
\usepackage{microtype}

% Helvetica Neue no trae la flecha U+2192, asi que cada «→» escrito literalmente
% en el YAML se compilaba a NADA, en silencio (462 flechas en 61 guias: rutas del
% tipo «paleta → tipografia → lockup» salian rotas). Mapeamos el caracter a su
% version matematica, que si tiene glifo. Asi el autor puede seguir escribiendo
% «→» comodamente en el YAML.
\usepackage{newunicodechar}
\newunicodechar{→}{\ensuremath{\rightarrow}}

\setmainfont{Helvetica Neue}
\setsansfont{Helvetica Neue}
\setlength{\parindent}{0pt}
\setlength{\parskip}{6pt}
\hyphenpenalty=200
\exhyphenpenalty=200
\pretolerance=400
\tolerance=2000
\setlength{\emergencystretch}{1em}
\renewcommand{\arraystretch}{1.25}
\setlist{leftmargin=5mm,itemsep=1mm,topsep=1mm}
\newcolumntype{Y}{>{\RaggedRight\arraybackslash}X}

\definecolor{milcMagenta}{HTML}{E6007E}
\definecolor{milcVino}{HTML}{5A0038}
\definecolor{milcTurquesa}{HTML}{0093A5}
\definecolor{milcVerde}{HTML}{58A923}
\definecolor{milcMostaza}{HTML}{E5B400}
\definecolor{milcOcre}{HTML}{B86600}
\definecolor{milcNegro}{HTML}{241816}
\definecolor{milcGris}{HTML}{F7F7F4}
\definecolor{milcLinea}{HTML}{E8E2DD}
\definecolor{milcGradePrimary}{HTML}{84CC16}
\definecolor{milcGradeDark}{HTML}{4D7C0F}
\definecolor{milcGradeSoft}{HTML}{ECFCCB}
\definecolor{milcGradeAccent}{HTML}{CFE7FF}
\definecolor{milcGradeText}{HTML}{FFFFFF}
\color{milcNegro}

\pagestyle{fancy}
\fancyhf{}
\lhead{\small Tecnología e Informática · Grado 7}
\rhead{\small Guía 11 · MILC v3}
\cfoot{\small\thepage}
\renewcommand{\headrulewidth}{0pt}

\newtcolorbox{titlebox}[1]{
  enhanced,colback=#1,colframe=#1,arc=7mm,boxrule=0pt,
  left=7mm,right=7mm,top=3mm,bottom=3mm,
  before skip=8pt,after skip=8pt,
  fontupper=\sffamily\bfseries\Large\color{white}
}

\newtcolorbox{softbox}[3]{
  enhanced,colback=#2,colframe=#1,borderline west={4pt}{0pt}{#1},
  arc=2mm,boxrule=.4pt,left=4mm,right=4mm,top=3mm,bottom=3mm,
  before skip=6pt,after skip=8pt,title={#3},coltitle=white,
  colbacktitle=#1,fonttitle=\sffamily\bfseries,fontupper=\RaggedRight,
  attach boxed title to top left={xshift=3mm,yshift=-2mm},
  boxed title style={arc=2mm, boxrule=0pt}
}

\newtcolorbox{drawbox}[2]{
  enhanced,colback=white,colframe=#1,arc=4mm,boxrule=1pt,
  left=5mm,right=5mm,top=4mm,bottom=4mm,before skip=8pt,after skip=8pt,
  title={#2},coltitle=white,colbacktitle=#1,fonttitle=\sffamily\bfseries,
  fontupper=\RaggedRight,
  attach boxed title to top left={xshift=4mm,yshift=-2mm},
  boxed title style={arc=3mm, boxrule=0pt}
}

\newtcolorbox{infoband}[1]{
  enhanced,colback=#1!8,colframe=#1!50,arc=2mm,boxrule=.5pt,
  left=4mm,right=4mm,top=2mm,bottom=2mm,before skip=4pt,after skip=4pt,
  fontupper=\small\RaggedRight
}

% Puente narrativo entre secciones (contrato editorial Conéctate).
% Da continuidad: "Acabas de X. Ahora vas a Y porque Z."
% Visualmente: banda izquierda en color del grado, texto en itálica pequeña.
\newtcolorbox{puentebox}{
  enhanced,colback=milcGradeSoft,colframe=milcGradeDark,
  borderline west={3pt}{0pt}{milcGradeDark},
  arc=0mm,boxrule=0pt,left=5mm,right=4mm,top=2.5mm,bottom=2.5mm,
  before skip=4pt,after skip=8pt,
  fontupper=\itshape\small\color{milcNegro}\RaggedRight
}
% La flecha va en modo matematico a proposito: Helvetica Neue NO tiene el glifo
% U+2192, asi que \textrightarrow se compilaba a NADA («Missing character: There
% is no → in font Helvetica Neue») y el puente salia sin flecha. En modo
% matematico la flecha viene de la fuente matematica y si se dibuja.
\newcommand{\puente}[1]{%
  \begin{puentebox}%
    \textcolor{milcGradeDark}{$\rightarrow$}\hspace{2mm}#1%
  \end{puentebox}%
}

\newcommand{\linead}[1]{\noindent\color{milcLinea}\rule{#1}{0.5pt}\color{milcNegro}}
\newcommand{\respuesta}[1]{\par\vspace{2mm}\foreach \n in {1,...,#1}{\noindent\color{milcLinea}\rule{\linewidth}{0.5pt}\par\vspace{5mm}}\color{milcNegro}}
\newcommand{\checkbox}{\textcolor{milcGradeDark}{\fbox{\phantom{x}}}\hspace{5pt}}

\newcommand{\phasecard}[4]{
\begin{tcolorbox}[enhanced,colback=#3,colframe=#2,arc=3mm,boxrule=.5pt,height=3.05cm,valign=center,halign=center,left=1.5mm,right=1.5mm,top=2mm,bottom=2mm]
{\sffamily\bfseries\fontsize{22}{24}\selectfont\textcolor{#2}{#1}}\par
{\sffamily\bfseries\small #4}
\end{tcolorbox}
}

% ── Recursos visuales de la guía ──────────────────────────────────────
% Declarados en el bloque `recursos:` del YAML y emitidos por el builder.
% \guiaFigura   → imagen rasterizada o PDF (foto, carta celeste, montaje).
% \guiaDiagrama → fuente TikZ/circuitikz (.tex) incrustada como vector:
%                 es la vía para circuitos y esquemáticos editables.
% #1 = ruta relativa al .tex · #2 = pie (puede ir vacío)
\newcommand{\guiaFigura}[2]{%
\begin{center}
\includegraphics[width=0.88\linewidth,height=0.38\textheight,keepaspectratio]{#1}\\[1.5mm]
{\footnotesize\itshape\textcolor{milcNegro!75}{#2}}
\end{center}
}

\newcommand{\guiaDiagrama}[2]{%
\begin{center}
\input{#1}\\[1.5mm]
{\footnotesize\itshape\textcolor{milcNegro!75}{#2}}
\end{center}
}

\begin{document}
\thispagestyle{empty}
\begin{tikzpicture}[remember picture,overlay]
  \fill[white] (current page.south west) rectangle (current page.north east);
  \path[fill=milcGradePrimary,rounded corners=16mm]
    ([xshift=.55cm,yshift=-.55cm]current page.north west) rectangle
    ([xshift=-.55cm,yshift=3.65cm]current page.south east);

  \node[anchor=north west,text=milcGradeText] at ([xshift=2.0cm,yshift=-1.85cm]current page.north west)
    {\sffamily\bfseries\fontsize{14}{16}\selectfont GRADO};
  \node[anchor=north west,text=milcGradeText,opacity=.82] at ([xshift=2.0cm,yshift=-2.75cm]current page.north west)
    {\sffamily\bfseries\fontsize{9}{11}\selectfont SERIE GUÍAS · TECNOLOGÍA E INFORMÁTICA};

  \begin{scope}[shift={([xshift=-3.05cm,yshift=-2.55cm]current page.north east)}]
    \fill[white,opacity=.80] (-.62,-.52) rectangle (.62,.52);
    \draw[milcGradeText,line width=1pt] (-.62,-.52) rectangle (.62,.52);
    \fill[milcGradeDark] (-.42,-.38) rectangle (-.22,.06);
    \fill[milcGradeAccent] (-.10,-.38) rectangle (.10,.24);
    \fill[milcGradePrimary!60!black] (.22,-.38) rectangle (.42,.38);
  \end{scope}

  \node[anchor=west,text=milcGradeText] at ([xshift=1.9cm,yshift=-12.85cm]current page.north west)
    {\sffamily\bfseries\fontsize{124}{124}\selectfont 7};
  % El circulito de grado (°) solo aplica a las guias de grado («11°»); en
  % Bebras y Semillero el numero grande es un momento/modulo, no un grado.
  \draw[milcGradeText,line width=1.7pt]
    ([xshift=4.52cm,yshift=-11.68cm]current page.north west) circle (.13cm);

  \node[anchor=west,text=milcGradeText,text width=8.7cm,align=left] at ([xshift=10.35cm,yshift=-11.6cm]current page.north west)
    {
     {\sffamily\bfseries\fontsize{19}{22}\selectfont Apertura\\periodo 2 ---\\pensar como\\teje la\\abuela}\\[4mm]
     {\sffamily\bfseries\fontsize{23}{26}\selectfont Guía 11-7-TIC}\\[2mm]
     {\sffamily\fontsize{13}{16}\selectfont Período 2 · Algoritmia y pensamiento computacional}\\[5mm]
     {\sffamily\bfseries\fontsize{11}{13}\selectfont Institución Educativa Sor María Juliana}\\[1mm]
     {\sffamily\fontsize{10}{12}\selectfont Autor: Dr. Álvaro Cárdenas Orozco}
    };

  \path[fill=milcGradeSoft,rounded corners=7mm]
    ([xshift=1.35cm,yshift=.85cm]current page.south west) rectangle
    ([xshift=-1.35cm,yshift=4.25cm]current page.south east);
  \draw[milcGradeDark,line width=.65pt,rounded corners=7mm]
    ([xshift=1.35cm,yshift=.85cm]current page.south west) rectangle
    ([xshift=-1.35cm,yshift=4.25cm]current page.south east);
  \node[anchor=center,text=milcNegro,text width=17.0cm,align=center] at ([yshift=2.55cm]current page.south)
    {\sffamily\fontsize{10}{12}\selectfont Metodología MILC · Apertura ancestral · Pensamiento computacional · Triángulo Dussel-Estoicismo-Floridi\\
    \textcolor{milcGradeDark}{\bfseries Producto:} Tu \textbf{cuaderno listo para el periodo 2 de grado 7}: encabezado fijo + \textbf{índice del P2} con las 10 sesiones + \textbf{compromiso firmado} de 5 acuerdos del pensamiento computacional (paciencia, paso a paso, revisar errores, intentar de nuevo, compartir aprendizajes) + \textbf{5 preguntas-radar} sobre algoritmia y Scratch + \textbf{mini-investigación} sobre el tejido tradicional del Cauca (canastos, mochilas, sombreros) y por qué tejer es \emph{algorítmico}.};
\end{tikzpicture}
\mbox{}
\newpage

\section*{Apertura: Apertura periodo 2 --- pensar como teje la abuela}

\begin{softbox}{milcGradeDark}{milcGradeSoft}{Saber ancestral}
\textbf{En el Valle del Cauca, en las veredas que rodean Cartago, las abuelas tejedoras sabían programar mucho antes de que se inventaran los computadores.} Cuando doña Sofía, una abuela de la vereda La Cabaña, se sentaba a tejer un canasto de fibra de plátano, no improvisaba: \textbf{seguía un algoritmo}. Empezaba con 8 hebras cruzadas en el centro (\emph{paso 1}), levantaba la hebra de la derecha, la pasaba sobre 2 hebras y por debajo de la siguiente (\emph{paso 2 que se repite}), contaba hasta llegar a 12 puntos (\emph{condicional: si llegó a 12, cambio de hebra}), y volvía a comenzar el patrón. Después de 1 hora, tenía las paredes del canasto. \textbf{Si te detenías a verla, te dabas cuenta:} doña Sofía hacía \emph{secuencia} (pasos en orden), \emph{repetición} (el patrón se repetía cientos de veces), \emph{condicionales} (cuando llegaba al borde, cambiaba), \emph{variables} (contaba los puntos en su cabeza). \textbf{Eso es exactamente lo que hace un computador.} Las mochilas Wayuu de La Guajira, los sombreros aguadeños del Quindío, los canastos del Cauca: todos son algoritmos visuales transmitidos durante siglos sin necesidad de pantallas. Este periodo te enseña a hacer lo mismo, pero en lenguaje computacional.
\end{softbox}

\begin{softbox}{milcTurquesa}{milcGris}{Saber contemporáneo}
Este periodo se llama \textbf{Algoritmia y pensamiento computacional}. ``Algoritmo'' es la palabra técnica para \emph{una receta paso a paso}. \textbf{Pensamiento computacional} es la habilidad de \emph{descomponer un problema complejo en pasos manejables que un computador (o cualquier persona) puede ejecutar}. Vas a aprenderlo en 10 sesiones: \textbf{(S1) Hoy: Apertura.} \textbf{(S2) ¿Qué es un algoritmo?} \textbf{(S3) Pensamiento computacional --- las 4 técnicas.} \textbf{(S4) Diagramas de flujo.} \textbf{(S5) Variables.} \textbf{(S6) Condicionales (si... entonces).} \textbf{(S7) Bucles (repeticiones).} \textbf{(S8) Scratch --- mi primer programa.} \textbf{(S9) Scratch --- narrativa interactiva.} \textbf{(S10) Cosecha --- mi primer programa completo.} \textbf{Scratch} es el lenguaje de programación visual del MIT (gratis, en el navegador) que vas a usar en las últimas 3 sesiones. Funciona con bloques que se encajan como Lego. \textbf{Al final del periodo, vas a saber programar tu primera historia interactiva}. Eso te abre la puerta a programar de verdad en grados 8°-11° y eventualmente en universidad si te animas.
\end{softbox}

\begin{softbox}{milcMostaza}{milcGris}{Pregunta-puente}
\textit{Si te pidieran explicarle a un robot \emph{``cómo hacer una arepa''}, ¿\textbf{podrías hacerlo}? ¿Cuántos pasos tendrías que decirle? ¿Qué pasaría si te olvidaras de un paso?}
\end{softbox}

\begin{softbox}{milcVerde}{milcGris}{Saber y saber hacer}
Al terminar la clase: (1) podrás \textbf{identificar} la ruta del periodo 2; (2) sabrás \textbf{explicar} por qué tejer es ``algorítmico''; (3) podrás \textbf{aplicar} los 5 acuerdos del pensamiento computacional; (4) habrás \textbf{creado} tu cuaderno P2 firmado + investigación.
\end{softbox}



\small
\begin{tabularx}{\textwidth}{>{\bfseries}p{3.0cm}Y}
\rowcolor{milcGradePrimary!12}Autor & Dr. Álvaro Cárdenas Orozco\\
DBA & Construye algoritmos y diagramas de flujo para resolver problemas paso a paso (MEN, T\&I 7°)\\
Referentes & Saber ancestral del tejedor · Pensamiento computacional · Scratch\\
Producto final & Tu \textbf{cuaderno listo para el periodo 2 de grado 7}: encabezado fijo + \textbf{índice del P2} con las 10 sesiones + \textbf{compromiso firmado} de 5 acuerdos del pensamiento computacional (paciencia, paso a paso, revisar errores, intentar de nuevo, compartir aprendizajes) + \textbf{5 preguntas-radar} sobre algoritmia y Scratch + \textbf{mini-investigación} sobre el tejido tradicional del Cauca (canastos, mochilas, sombreros) y por qué tejer es \emph{algorítmico}.\\
\end{tabularx}
\normalsize

\puente{Hoy haces 4 cosas: descubres por qué tejer es programar, conoces la ruta, firmas el compromiso, investigas el tejido tradicional.}

\begin{titlebox}{milcGradeDark}
Ruta MILC: así vamos a aprender
\end{titlebox}
\begin{tabularx}{\textwidth}{>{\centering\arraybackslash}X>{\centering\arraybackslash}X>{\centering\arraybackslash}X>{\centering\arraybackslash}X}
\phasecard{1}{milcVerde}{milcGris}{Escucha\\[-1mm]\small Observo, escucho y pregunto.} &
\phasecard{2}{milcTurquesa}{milcGris}{Sistematización\\[-1mm]\small Pensamiento computacional.} &
\phasecard{3}{milcMagenta}{milcGris}{Praxis\\[-1mm]\small Construyo y aplico.} &
\phasecard{4}{milcVino}{milcGris}{Evaluación\\[-1mm]\small Reflexión liberadora.}
\end{tabularx}


\puente{Empezamos imaginando a la tejedora paso a paso.}

\begin{titlebox}{milcVerde}
Fase 1 · Escucha
\end{titlebox}

\begin{softbox}{milcVerde}{milcGris}{Escena para escuchar}
\textbf{Actividad 1 · IDENTIFICA --- Imagina a doña Sofía tejiendo} (12 min · individual). Cierra los ojos e imagina la escena. Doña Sofía está sentada en una silla baja, con un canasto a medio tejer entre las manos. Ella sigue \textbf{el algoritmo del tejido}. En el cuaderno, responde con tu intuición: \emph{(1) ¿Cuántos pasos diferentes crees que tiene el algoritmo del tejido? (estima)}. \emph{(2) ¿Cuáles pasos crees que se REPITEN varias veces?}. \emph{(3) ¿Qué pasaría si doña Sofía olvida un paso?}. \emph{(4) ¿Cómo sabe doña Sofía cuándo terminar (cuándo el canasto está completo)?}. \emph{(5) ¿Doña Sofía aprendió este algoritmo en una clase, o de otra manera?}. Tus respuestas no son ``correctas'' o ``incorrectas''; son tu intuición inicial.
\end{softbox}

\checkbox Cerré los ojos e imaginé la escena.\\
\checkbox Respondí las 5 preguntas con honestidad.\\
\checkbox Pensé sin copiar al compañero.

\begin{infoband}{milcVerde}


\textbf{Lo que va al cuaderno (Actividad 1 · IDENTIFICA).} Título: «Actividad 1 --- Imagino a la tejedora». \textbf{Formato:} 5 respuestas cortas. \textbf{Extensión:} 5 líneas. \textbf{Sabes que terminaste cuando} reconoces que el tejido tiene estructura algorítmica.
\end{infoband}

\puente{Después aprendes los 5 acuerdos del pensamiento computacional + ruta del P2.}

\begin{titlebox}{milcTurquesa}
Fase 2 · Sistematización con pensamiento computacional
\end{titlebox}

\begin{softbox}{milcTurquesa}{milcGris}{Cuatro pilares aplicados al tema}
Lo que descubriste imaginando a doña Sofía es \textbf{pensamiento computacional}. Ahora aprende la ruta del periodo + los 5 acuerdos para llegar bien.
\end{softbox}

\small
\begin{tabularx}{\textwidth}{>{\bfseries\RaggedRight\arraybackslash}p{3.6cm}Y}
\rowcolor{milcTurquesa!90}\color{white}Pilar & \color{white}\bfseries Aplicación al tema\\
1. Descomposición & \textbf{La ruta del P2 --- 10 sesiones organizadas en 3 bloques.} \textbf{Bloque 1 (S2-S4): Algoritmo y diagramas.} Qué es un algoritmo, las 4 técnicas del pensamiento computacional, cómo dibujar un algoritmo. \textbf{Bloque 2 (S5-S7): Componentes de programación.} Variables, condicionales, bucles --- los 3 ladrillos básicos de cualquier programa. \textbf{Bloque 3 (S8-S10): Scratch.} El lenguaje visual del MIT donde aplicas todo lo anterior con bloques. Al final, tu \emph{primer programa interactivo completo}.\\
2. Reconocimiento de patrones & \textbf{¿Qué es Scratch?} Lenguaje de programación visual creado en el MIT (Estados Unidos) para que niños y jóvenes aprendan a programar. \textbf{No requiere escribir código:} los \emph{bloques} de instrucciones se arrastran y se encajan como piezas de Lego. \textbf{Funciona en cualquier computador con navegador:} entras a \texttt{scratch.mit.edu}, creas tu cuenta gratis (o usas la cuenta institucional), y empiezas a programar. \textbf{Más de 100 millones de proyectos} subidos por chicos de todo el mundo. Algunos son juegos, otros cuentos, otros simuladores. Tú harás los tuyos al final del periodo.\\
3. Abstracción & \textbf{Los 5 acuerdos del pensamiento computacional.} (1) \textbf{Paciencia:} programar requiere pensar despacio. La primera vez tarda; con práctica se vuelve rápido. (2) \textbf{Paso a paso:} no querer hacer todo de una. Un paso, después otro, después otro. (3) \textbf{Revisar errores con calma:} cuando algo no funciona, eso no es fracaso --- es información. Te dice qué corregir. \emph{Los programadores experimentados pasan más tiempo revisando errores que escribiendo código nuevo}. (4) \textbf{Intentar de nuevo:} si la primera versión no funciona, modificas y vuelves a intentar. \emph{Ningún programa funciona perfecto en el primer intento}. (5) \textbf{Compartir aprendizajes:} cuando descubres algo, lo compartes con el equipo. Programar es deporte colectivo, no solitario.\\
4. Algoritmo conceptual & \textbf{4 errores típicos al iniciar pensamiento computacional.} (1) \emph{``No soy bueno para esto''} antes de intentar --- nadie nace programador, todos aprenden. (2) \emph{Querer hacer un programa complicado el primer día} --- la S8 empieza con \emph{``el sprite que se mueve''}, no con un videojuego. (3) \emph{Saltarse pasos del razonamiento} --- en algoritmia, saltar 1 paso a veces invalida todo. (4) \emph{Frustrarse cuando aparece un error} --- los errores son tu maestro principal. Sin error, no aprendes a depurar.\\
\end{tabularx}
\normalsize

\begin{drawbox}{milcTurquesa}{Ruta P2 + 5 acuerdos}
\textbf{Ruta P2:} \\[1mm]
Bloque 1 (S2-S4): algoritmo + 4 técnicas + diagramas. \\[1mm]
Bloque 2 (S5-S7): variables + condicionales + bucles. \\[1mm]
Bloque 3 (S8-S10): Scratch + narrativa + cosecha. \\[1mm]
\textbf{5 acuerdos:} paciencia · paso a paso · revisar errores · intentar de nuevo · compartir.
\end{drawbox}

\begin{softbox}{milcMostaza}{milcGris}{Errores comunes y su corrección}
\textbf{Mitos sobre la programación que vas a oír.} (1) \emph{``Solo los genios programan''} --- Falso. Programar es habilidad que se aprende, como leer. (2) \emph{``Hay que ser bueno en matemáticas para programar''} --- Falso. La mate ayuda pero no es requisito. La habilidad central es \emph{pensar paso a paso}, no resolver ecuaciones. (3) \emph{``Programar es solo para hombres''} --- Totalmente falso. Las primeras programadoras fueron mujeres (Ada Lovelace, las ``chicas del ENIAC''). (4) \emph{``Programar en Scratch no es programar de verdad''} --- Falso. Scratch usa los mismos conceptos (variables, condicionales, bucles, eventos) que cualquier lenguaje profesional.
\end{softbox}

\begin{infoband}{milcTurquesa}


\textbf{Lo que va al cuaderno (Actividad 2 · EXPLICA).} Título: «Actividad 2 --- Ruta P2 + 5 acuerdos». \textbf{Formato:} bloque A con los 3 bloques del periodo (lista). Bloque B con los 5 acuerdos firmados. \textbf{Extensión:} 12 líneas. \textbf{Sabes que terminaste cuando} podrías recitar los 5 acuerdos sin abrir el cuaderno.
\end{infoband}

\puente{Con todo claro, dejas el cuaderno listo + 5 preguntas + investigación.}

\begin{titlebox}{milcMagenta}
Fase 3 · Praxis con pensamiento computacional
\end{titlebox}

\begin{softbox}{milcMagenta}{milcGris}{Construyendo el producto con los cuatro pilares}
\textbf{Actividad 3 · CREA --- Cuaderno P2 + 5 preguntas + investigación del tejido} (28 min · individual + tarea en casa). Vas a dejar el cuaderno listo + plantear las preguntas que llevas al periodo + investigar el tejido tradicional como tarea para casa.
\end{softbox}

\small
\begin{tabularx}{\textwidth}{>{\bfseries\RaggedRight\arraybackslash}p{3.6cm}Y}
\rowcolor{milcMagenta!90}\color{white}Pilar & \color{white}\bfseries Aplicación al producto\\
Descomposición & \textbf{Paso 1 --- Índice del P2 + encabezado fijo.} En tu cuaderno de Tecnología (el mismo de P1), abre una página nueva. Escribe el título: \textbf{``Periodo 2 --- Algoritmia y pensamiento computacional''}. Debajo, índice con las 10 sesiones del P2 (mira la ruta de la sistematización).\\
Patrones & \textbf{Paso 2 --- Los 5 acuerdos firmados.} Después del índice, lista numerada de los 5 acuerdos (paciencia, paso a paso, revisar errores, intentar de nuevo, compartir). Al lado de cada uno, palomita ✓ si te comprometes. Si dudas con alguno, conversa con el profe.\\
Abstracción & \textbf{Paso 3 --- 5 preguntas-radar.} Después de los acuerdos, lista numerada de 5 preguntas tuyas que esperas que el P2 responda. Sugerencias para inspirarte: \emph{``¿Por qué dicen que tejer es como programar?''}, \emph{``¿Es difícil programar en Scratch?''}, \emph{``¿Puedo hacer un videojuego al final del periodo?''}, \emph{``¿Cómo se llaman los errores en programación?''}.\\
Algoritmo de armado & \textbf{Paso 4 --- Mini-investigación del tejido.} En la siguiente hoja del cuaderno, abre el espacio: \textbf{``Mini-investigación: el tejido como algoritmo''}. Esta investigación la haces \emph{en casa} esta semana (no en clase). Tareas: \emph{(a) Encuentra a una persona mayor de tu familia o barrio que sepa tejer (puede ser canasto, manilla, suéter, mochila, sombrero). (b) Pregúntale: ¿cuántos pasos básicos tiene el tejido?, ¿qué pasos se repiten?, ¿cómo aprendió ella a tejer?, ¿qué hace si se equivoca?. (c) Anota mínimo 5 datos concretos en el cuaderno.}. Si no tienes a nadie en casa, busca en internet con criterio (aplicando CRAAP de G6).\\
\end{tabularx}
\normalsize

\begin{drawbox}{milcMagenta}{Antes de cerrar la S1 del P2}
\checkbox El índice del P2 está con 10 sesiones. \\[1mm]
\checkbox Los 5 acuerdos están firmados con palomita. \\[1mm]
\checkbox Tengo 5 preguntas-radar reales. \\[1mm]
\checkbox Tengo plan para la mini-investigación del tejido. \\[1mm]
\checkbox La sesión 1 tiene encabezado fijo y registro. \\[1mm]
\checkbox Está firmado con nombre y fecha.
\end{drawbox}

\begin{softbox}{milcMagenta}{milcGris}{Plantilla de trabajo}
\textbf{Estructura.} \textit{Hoja 1:} título P2 + índice de 10 sesiones. \textit{Hoja 2:} 5 acuerdos firmados + 5 preguntas-radar. \textit{Hoja 3:} mini-investigación (a completar en casa). \textit{Cuaderno principal:} S1 con encabezado fijo.
\end{softbox}

\begin{infoband}{milcMagenta}


\textbf{Lo que va al cuaderno (Actividad 3 · CREA).} Título: «Actividad 3 --- Cuaderno inaugural G7-P2». \textbf{Formato:} 3 hojas iniciales + S1 registrada. \textbf{Extensión:} 3 hojas + sesión. \textbf{Sabes que terminaste cuando} tu cuaderno está listo para empezar la S2 sin reorganizar nada.
\end{infoband}

\puente{El producto es el cuaderno P2 inaugural + mini-investigación + compromiso firmado.}

\begin{titlebox}{milcMostaza}
Producto de la sesión
\end{titlebox}
\textbf{Cuaderno P2 inaugural + investigación del tejido + 5 acuerdos firmados}.
\begin{softbox}{milcMostaza}{milcGris}{Criterios de calidad}
(1) Índice del P2 con 10 sesiones. (2) 5 acuerdos firmados con palomita. (3) 5 preguntas-radar personales y específicas. (4) Plan claro de mini-investigación del tejido (con persona o fuente identificada). (5) S1 con encabezado fijo y registro de actividades.
\end{softbox}

\puente{Antes de salir, verifica que tu cuaderno tiene los 5 elementos del periodo.}

\begin{titlebox}{milcVino}
Fase 4 · Evaluación liberadora — cinco dimensiones
\end{titlebox}

\begin{softbox}{milcVino}{milcGris}{Sentido de la evaluación}
Evaluar no es solo revisar una nota. Es reconocer qué aprendí, qué cambió en mí, qué emociones manejé, cómo me relaciono con la ciudad y la comunidad, y qué le dejaré a quien viene detrás.
\end{softbox}

\textbf{Mi reflexión liberadora en cinco dimensiones.} Completa cada frase iniciada:

\small
\begin{tabularx}{\textwidth}{>{\bfseries\RaggedRight\arraybackslash}p{3.4cm}Y>{\RaggedRight\arraybackslash}p{4.6cm}}
\rowcolor{milcVino}\color{white}Dimensión & \color{white}\bfseries Mi respuesta (frame starter) & \color{white}\bfseries Algo que descubrí\\
1. Personal & Lo que más cambió en mí fue \linead{6.5cm} & \linead{4.2cm}\\
2. Emocional & La emoción más fuerte fue \linead{2.5cm} y la regulé \linead{3cm} & \linead{4.2cm}\\
3. Ciudadana & En democracia esto importa porque \linead{6cm} & \linead{4.2cm}\\
4. Local (Cartago) & En mi barrio sería útil para \linead{6.5cm} & \linead{4.2cm}\\
5. Intergeneracional & A mi abuela/abuelo le diría \linead{6.5cm} & \linead{4.2cm}\\
\end{tabularx}
\normalsize

\begin{infoband}{milcVino}
\textbf{Para la próxima sustentación.} Vuelve a esta tabla en seis meses. Verás que las respuestas cambian — no porque lo aprendido cambie, sino porque tú cambias. La evaluación liberadora es una conversación contigo a través del tiempo.
\end{infoband}


\puente{Cerramos con 3 voces sobre por qué pensar paso a paso es de las habilidades más poderosas que existen.}

\begin{titlebox}{milcGradeDark}
Triángulo de pensamiento · cierre filosófico
\end{titlebox}

\begin{softbox}{milcVino}{milcGris}{Dussel · lente del nosotros}
\textit{````Las abuelas tejedoras eran programadoras antes de que se inventara el computador. El saber popular precede a la academia.''''}

La academia muchas veces presenta el pensamiento computacional como algo \emph{moderno e importado}. Pero \textbf{las abuelas tejedoras Wayuu, las cesteras del Cauca, las artesanas del Pacífico} lo dominan desde hace siglos. Cuando tejen un patrón complejo, están ejecutando un algoritmo mental: \emph{secuencia, repetición, condicional, variable}. Reconocer eso le devuelve dignidad al saber popular y te recuerda que \emph{no estás aprendiendo algo ajeno; estás aprendiendo en lenguaje computacional algo que tu cultura ya conoce}.

\textbf{¿En mi familia hay alguien que ``programa'' sin saberlo (cocina con receta exacta, teje, hace algo paso a paso)? ¿Qué les podría enseñar el pensamiento computacional explícito?}
\end{softbox}
\linead{\linewidth}\\[1mm]
\linead{\linewidth}

\begin{softbox}{milcMostaza}{milcGris}{Estoicismo · Marco Aurelio · cuidado interior}
\textit{````Quien aprende a pensar paso a paso domina cualquier oficio. La paciencia con el orden es el secreto del trabajo bien hecho.''''}

\textbf{La habilidad central del pensamiento computacional es la paciencia paso a paso}. Ni los más rápidos pueden programar sin desglosar el problema en pasos. Quien se apresura, fracasa. Quien va paso a paso, llega. \emph{Esa disciplina --- la paciencia ordenada --- es útil mucho más allá de la programación}: para resolver problemas de matemáticas, organizar viajes, planear estudios, criar hijos. Aprender pensamiento computacional es aprender una forma adulta de pensar.

\textbf{¿Tiendo a apurarme cuando un problema es difícil, o sé descomponer paso a paso?}
\end{softbox}
\linead{\linewidth}\\[1mm]
\linead{\linewidth}

\begin{softbox}{milcTurquesa}{milcGris}{Floridi · lente de la infoesfera}
\textit{````Programar es la nueva alfabetización. No es para todos hacer carrera de programador, pero todos deben entender cómo se programa el mundo en que viven.''''}

Vivimos en un mundo \textbf{programado}: TikTok te recomienda videos por algoritmo, el banco aprueba créditos por algoritmo, los semáforos cambian por algoritmo. \emph{Si no entiendes nada de cómo funcionan los algoritmos, te conviertes en sujeto pasivo del mundo digital}. Entenderlos --- aunque no programes profesionalmente --- te hace ciudadano alfabetizado del siglo XXI. Este periodo es ese entrenamiento mínimo.

\textbf{¿Quiero ser sujeto pasivo de los algoritmos que me rodean, o ciudadano que los entiende y decide con criterio?}
\end{softbox}
\linead{\linewidth}\\[1mm]
\linead{\linewidth}

\begin{titlebox}{milcVino}
Compromiso final
\end{titlebox}

\small
\begin{tabularx}{\textwidth}{>{\RaggedRight\arraybackslash}p{3.0cm}YYY}
\rowcolor{milcVino}\color{white}\bfseries Criterio & \color{white}\bfseries Lo logré & \color{white}\bfseries En proceso & \color{white}\bfseries Necesito apoyo\\
Comprensión del tema & Explico ideas centrales con mis palabras. & Reconozco ideas pero falta precisión. & Necesito repasar conceptos base.\\
Pensamiento computacional & Apliqué los 4 pilares al diseñar mi producto. & Apliqué algunos pilares. & Necesito releer la fase 2.\\
Aplicación y producto & Mi producto cumple los criterios y se puede revisar. & Producto funcional con ajustes pendientes. & Aún en construcción.\\
Reflexión 5 dimensiones & Respondí cada dimensión con honestidad. & Respondí algunas con detalle. & Mis respuestas son muy generales.\\
Triángulo de pensamiento & Conecté las 3 voces con mi trabajo real. & Conecté al menos una. & No relaciono las voces con mi proceso.\\
\end{tabularx}
\normalsize

\textbf{Hoy entendí que:}\\[1mm]
\linead{\linewidth}\\[1mm]
\linead{\linewidth}

\textbf{Mi compromiso para la próxima sesión es:}\\[1mm]
\linead{\linewidth}\\[1mm]
\linead{\linewidth}

\begin{softbox}{milcMostaza}{milcGris}{Cita de despedida · Marco Aurelio}
\textit{``El obstáculo en el camino se vuelve el camino. Lo que estorba al camino se vuelve camino.''}\\[1mm]
Cada error de hoy, bien observado, se convierte en pieza del camino que pisarás mañana.
\end{softbox}

\small
\begin{tabularx}{\textwidth}{>{\bfseries}p{3.6cm}Y}
\rowcolor{milcGradeDark!12}Nombre completo: & \linead{10cm}\\
Fecha: & \linead{4cm} \quad Curso/grupo: \linead{4cm}\\
Mi firma: & \linead{9cm}\\
\end{tabularx}
\normalsize

\begin{infoband}{milcGradeDark}
\textbf{Para la próxima sesión.} Esta semana, completa la mini-investigación del tejido: habla con una persona mayor, anota 5 datos. La próxima clase es S2: \emph{¿Qué es un algoritmo?}. Vas a aprender la definición técnica + las 5 características de un buen algoritmo. Trae tu investigación: la usaremos.
\end{infoband}

\end{document}
